% ═══════════════════════════════════════════════════════════════════════════
% PLATZKARTEN: Kräfte zusammensetzen (UE 2)
% Physik Klasse 7 – 24 Karten auf 2 A4-Seiten
% ═══════════════════════════════════════════════════════════════════════════

\documentclass[a4paper,11pt]{article}

\usepackage[utf8]{inputenc}
\usepackage[T1]{fontenc}
\usepackage{helvet}
\renewcommand{\familydefault}{\sfdefault}

\usepackage[margin=10mm]{geometry}
\usepackage{tikz}
\usepackage{xcolor}
\usepackage{qrcode}

% Farbe
\definecolor{physik}{HTML}{2c5aa0}

% Konfiguration
\newcommand{\slurl}{https://devbox42.github.io/sciencesim/physik/kl7/kraft/02-kraft-groesse/MT.html}

\pagestyle{empty}

\newcommand{\platzkarte}[1]{%
\begin{tikzpicture}
    % Rahmen
    \draw[physik, line width=1.5pt] (0,0) rectangle (45mm,88mm);

    % Platznummer
    \node[anchor=north, physik, font=\fontsize{28}{32}\selectfont\bfseries]
        at (22.5mm,84mm) {P#1};

    % Namenslinie
    \draw[physik, line width=0.8pt] (5mm,58mm) -- (40mm,58mm);
    \node[anchor=north, font=\scriptsize, text=gray] at (22.5mm,57mm) {Name};

    % Niveau-Feld
    \draw[physik, line width=0.5pt] (5mm,38mm) rectangle (40mm,51mm);
    \node[anchor=north, font=\small\bfseries, physik] at (22.5mm,50mm) {Niveau};
    \node[anchor=center, font=\normalsize] at (22.5mm,43mm) {* \quad ** \quad ***};

    % QR-Code
    \node[anchor=south] at (22.5mm,2mm) {\qrcode[height=18mm]{\slurl}};
\end{tikzpicture}%
}

\begin{document}

% Seite 1: P1-P12
\noindent
\begin{tabular}{@{}c@{\hspace{2mm}}c@{\hspace{2mm}}c@{\hspace{2mm}}c@{}}
\platzkarte{1} & \platzkarte{2} & \platzkarte{3} & \platzkarte{4} \\[2mm]
\platzkarte{5} & \platzkarte{6} & \platzkarte{7} & \platzkarte{8} \\[2mm]
\platzkarte{9} & \platzkarte{10} & \platzkarte{11} & \platzkarte{12} \\
\end{tabular}

\newpage

% Seite 2: P13-P24
\noindent
\begin{tabular}{@{}c@{\hspace{2mm}}c@{\hspace{2mm}}c@{\hspace{2mm}}c@{}}
\platzkarte{13} & \platzkarte{14} & \platzkarte{15} & \platzkarte{16} \\[2mm]
\platzkarte{17} & \platzkarte{18} & \platzkarte{19} & \platzkarte{20} \\[2mm]
\platzkarte{21} & \platzkarte{22} & \platzkarte{23} & \platzkarte{24} \\
\end{tabular}

\end{document}
